% !TEX root = ../../../main.tex
% 来源：高中物理甲种本·第二册（热学·电学）
% 章节：分子运动论基础 / 物体是由分子组成的
% 原始文件：1.tex，行 206-231
% 类型：tikz
% ---
    \begin{tikzpicture}[>=latex, scale=.8]
\draw[dashed] (0,2) arc (90:150:2);
\draw[dashed] (0,2) arc (90:50:2);
\draw (0,6) arc (90:140:6);
\draw (0,6) arc (90:70:6);
\draw (0,0) [fill=black] circle (2.5pt)node[below]{$O$};
\draw (90:2) [fill=black] circle (2.5pt)node[right]{$b$};
\draw (90:6) [fill=black]circle (2.5pt)node[above]{$B$};
\draw (110:2) [fill=black]circle (2.5pt)node[left]{$a$};
\draw (110:6) [fill=black]circle (2.5pt)node[above]{$A$};

\draw[->] (0,0)--node[below]{$R$}(130:6);
\draw[very thick]  (0,0)--(90:6);
\draw[very thick]    (0,0)--(110:6);
\draw[very thick,->] (0,0)--(90:3.8);
\draw[very thick, ->] (0,0)--(110:3.8);
\draw [<->] (0,3) arc (90:110:3)node[right]{$\alpha$};
\draw[->] (0,0)--node[above]{$r$}(60:2);

\draw (1,4)--(90:4);
\draw (1,4)node[right]{离子的轨迹}--(110:3.3);

\draw (90:2)--(1,3)node[right]{钨原子};
\draw (110:2)--(1,3);

    \end{tikzpicture}
